\section{Domino tilings of height-3 rectangles}

\subsection{The tilings $A_n$, $B_n$, and $C$}

\begin{definition}
\label{def.gf.weighted-set.domino.Rn3.ABC}
\lean{AlgebraicCombinatorics.DominoTilings.TilingA, AlgebraicCombinatorics.DominoTilings.TilingB, AlgebraicCombinatorics.DominoTilings.TilingC}
\leantarget
\leanok
\textbf{(a)} For each even positive
integer $n$, we let $A_{n}$ be the domino tiling of $R_{n,3}$
consisting of the following dominos:

\begin{itemize}
\item the horizontal dominos $\left\{  \left(  2i-1,\ 1\right)  ,\ \left(
2i,\ 1\right)  \right\}  $ for all $i\in\left[  n/2\right]  $, which fill the
bottom row of $R_{n,3}$, and which we call the \textbf{basement dominos};

\item the vertical domino $\left\{  \left(  1,2\right)  ,\ \left(  1,3\right)
\right\}  $ in the first column, which we call the \textbf{left wall};

\item the vertical domino $\left\{  \left(  n,2\right)  ,\ \left(  n,3\right)
\right\}  $ in the last column, which we call the \textbf{right wall};

\item the horizontal dominos $\left\{  \left(  2i,\ 2\right)  ,\ \left(
2i+1,\ 2\right)  \right\}  $ for all $i\in\left[  n/2-1\right]  $, which fill
the middle row of $R_{n,3}$ (except for the first and last columns), and which
we call the \textbf{middle dominos};

\item the horizontal dominos $\left\{  \left(  2i,\ 3\right)  ,\ \left(
2i+1,\ 3\right)  \right\}  $ for all $i\in\left[  n/2-1\right]  $, which fill
the top row of $R_{n,3}$ (except for the first and last columns), and which we
call the \textbf{top dominos}.
\end{itemize}

\textbf{(b)} For each even positive integer $n$, we let $B_{n}$ be the domino
tiling of $R_{n,3}$ obtained by reflecting $A_{n}$ across the horizontal axis of
symmetry of $R_{n,3}$ (swapping row $1$ with row $3$ and fixing row $2$). \medskip

\textbf{(c)} We let $C$ denote the domino tiling of $R_{2,3}$ consisting of
three horizontal dominos:
$\left\{(1,1),(2,1)\right\}$,
$\left\{(1,2),(2,2)\right\}$, and
$\left\{(1,3),(2,3)\right\}$.
\end{definition}

\subsection{Properties of $A_n$, $B_n$, and $C$}

\begin{lemma}
\label{lem.details.domino.TilingA-hasTopVertical}
\lean{AlgebraicCombinatorics.DominoTilings.TilingA_hasTopVertical}
\leanhelper
\leanok
For each even $n \ge 2$, the tiling $A_n$ has a vertical domino in the top
two squares of column~$1$, namely the left wall $\{(1,2),(1,3)\}$.
\end{lemma}

\begin{proof}
\leanok
The left wall $\{(1,2),(1,3)\}$ is explicitly among the dominos of $A_n$,
is vertical, lies in column~$1$, and covers rows $2$--$3$.
\end{proof}

\begin{lemma}
\label{lem.details.domino.TilingB-hasBottomVertical}
\lean{AlgebraicCombinatorics.DominoTilings.TilingB_hasBottomVertical}
\leanhelper
\leanok
For each even $n \ge 2$, the tiling $B_n$ has a vertical domino in the bottom
two squares of column~$1$, namely the left wall $\{(1,1),(1,2)\}$ of $B_n$.
\end{lemma}

\begin{proof}
\leanok
The left wall of $B_n$ is $\{(1,1),(1,2)\}$, which is vertical, lies in
column~$1$, and covers rows $1$--$2$.
\end{proof}

\begin{lemma}
\label{lem.details.domino.TilingC-noVertical}
\lean{AlgebraicCombinatorics.DominoTilings.TilingC_noVerticalInCol1}
\leanhelper
\leanok
The tiling $C$ has no vertical domino in column~$1$. All three dominos
of $C$ are horizontal.
\end{lemma}

\begin{proof}
\leanok
$C$ consists of three horizontal dominos $\{(1,y),(2,y)\}$ for $y = 1, 2, 3$.
None is vertical.
\end{proof}

\subsection{Faultfreeness of $A_n$, $B_n$, and $C$}

\begin{lemma}
\label{lem.details.domino.TilingA-faultfree}
\lean{AlgebraicCombinatorics.DominoTilings.TilingA_isFaultfree}
\leanhelper
\leanok
For each even $n \ge 2$, the tiling $A_n$ is faultfree.
The basement dominos prevent faults between columns $2i-1$ and $2i$,
while the middle and top dominos prevent faults between columns $2i$ and $2i+1$.
\end{lemma}

\begin{proof}
\leanok
Let $k$ be a column with $1 \le k < n$.
If $k$ is even, say $k = 2j$, then the middle domino
$\{(2j, 2), (2j+1, 2)\}$ spans columns $2j$ and $2j+1$,
so there is no fault at column $k$.
If $k$ is odd, say $k = 2j+1$, then the basement domino
$\{(2j+1, 1), (2j+2, 1)\}$ spans columns $2j+1$ and $2j+2$,
so there is no fault at column $k$.
\end{proof}

\begin{lemma}
\label{lem.details.domino.TilingB-faultfree}
\lean{AlgebraicCombinatorics.DominoTilings.TilingB_isFaultfree}
\leanhelper
\leanok
For each even $n \ge 2$, the tiling $B_n$ is faultfree.
\end{lemma}

\begin{proof}
\leanok
Since $B_n$ is the reflection of $A_n$
(Lemma~\ref{lem.details.domino.TilingB-eq-reflect}), and reflection preserves
faultfreeness (Lemma~\ref{lem.details.domino.reflectTiling3-faultfree}),
$B_n$ is faultfree because $A_n$ is (Lemma~\ref{lem.details.domino.TilingA-faultfree}).
\end{proof}

\begin{lemma}
\label{lem.details.domino.TilingC-faultfree}
\lean{AlgebraicCombinatorics.DominoTilings.TilingC_isFaultfree}
\leanhelper
\leanok
The tiling $C$ is faultfree.
\end{lemma}

\begin{proof}
\leanok
The only potential fault in $R_{2,3}$ is at column $1$. But the domino
$\{(1,1),(2,1)\}$ spans columns $1$ and $2$, so there is no fault.
\end{proof}

\subsection{Reflection lemmas}

\begin{definition}
\label{def.details.domino.reflectTiling3}
\lean{AlgebraicCombinatorics.DominoTilings.reflectTiling3}
\leanhelper
\leanok
Given a tiling $T$ of $R_{n,3}$, the \emph{reflection} of $T$ is the tiling
obtained by applying the cell map $(x, y) \mapsto (x, 4 - y)$ to every cell
of every domino. This swaps rows $1 \leftrightarrow 3$ and fixes row $2$.
\end{definition}

\begin{proof}
\leanok
The reflected dominos are pairwise disjoint (since the cell map is injective)
and their union covers $R_{n,3}$ (since the cell map is a bijection on
$R_{n,3}$).
\end{proof}

\begin{lemma}
\label{lem.details.domino.reflectTiling3-minCol-maxCol}
\lean{AlgebraicCombinatorics.DominoTilings.reflectTiling3_dominos_minCol_maxCol}
\leanhelper
\leanok
The reflection of a tiling preserves the minimum and maximum column
of each domino: for every domino $d$ in $T$ and its reflected image $d'$,
the minimum column of $d'$ equals that of $d$, and the maximum column
of $d'$ equals that of $d$.
\end{lemma}

\begin{proof}
\leanok
The cell map $(x, y) \mapsto (x, 4-y)$ does not change the $x$-coordinate.
Hence $\min$ and $\max$ of the $x$-coordinates of the cells are preserved.
\end{proof}

\begin{lemma}
\label{lem.details.domino.reflectTiling3-faultfree}
\lean{AlgebraicCombinatorics.DominoTilings.reflectTiling3_isFaultfree}
\leanhelper
\leanok
If $T$ is a faultfree tiling of $R_{n,3}$, then its horizontal reflection
is also faultfree.
\end{lemma}

\begin{proof}
\leanok
Reflection only changes row coordinates, not column coordinates.
Therefore the minimum and maximum column of each domino are
preserved (Lemma~\ref{lem.details.domino.reflectTiling3-minCol-maxCol}),
and faults (which depend only on column spans) are preserved.
\end{proof}

\begin{lemma}
\label{lem.details.domino.reflectTiling3-top-bottom}
\lean{AlgebraicCombinatorics.DominoTilings.reflectTiling3_hasTopVertical_iff_hasBottomVertical}
\leanhelper
\leanok
A tiling $T$ of $R_{n,3}$ has a vertical domino in the top two squares
of column $c$ if and only if its reflection has a vertical domino in the
bottom two squares of column $c$ (and vice versa).
\end{lemma}

\begin{proof}
\leanok
The reflection swaps rows $1 \leftrightarrow 3$ and fixes row $2$.
A vertical domino covering rows $2$--$3$ in column $c$ is sent to one
covering rows $1$--$2$ in column $c$, and vice versa.
\end{proof}

\begin{lemma}
\label{lem.details.domino.reflectTiling3-involutive}
\lean{AlgebraicCombinatorics.DominoTilings.reflectTiling3_involutive}
\leanhelper
\leanok
Reflecting a tiling of $R_{n,3}$ twice yields the original tiling.
\end{lemma}

\begin{proof}
\leanok
The underlying cell reflection $(x, y) \mapsto (x, 4 - y)$ is an involution
on cells with $1 \le y \le 3$.
\end{proof}

\begin{lemma}
\label{lem.details.domino.TilingB-eq-reflect}
\lean{AlgebraicCombinatorics.DominoTilings.TilingB_eq_reflectTiling3_TilingA}
\leanhelper
\leanok
For each even $n \ge 2$, $B_n$ equals the reflection of $A_n$.
\end{lemma}

\begin{proof}
\leanok
By direct computation of the reflected domino sets.
The basement dominos of $A_n$ (in row $1$) reflect to the ``basement'' dominos
of $B_n$ (in row $3$), the left wall $\{(1,2),(1,3)\}$ reflects to
$\{(1,1),(1,2)\}$, and similarly for the right wall, middle, and top/bottom dominos.
\end{proof}

\begin{lemma}
\label{lem.details.domino.reflectTiling3-preserves-equiv}
\lean{AlgebraicCombinatorics.DominoTilings.reflectTiling3_preserves_TilingEquiv}
\leanhelper
\leanok
If two tilings $T_1$ and $T_2$ of $R_{n,3}$ are equal,
then their reflections are also equal.
\end{lemma}

\begin{proof}
\leanok
Reflection acts on the cells of each domino, so if $T_1$ and $T_2$ are
equal, so are their reflections.
\end{proof}

\subsection{Auxiliary lemmas for the classification}

\begin{lemma}
\label{lem.details.domino.parity}
\lean{AlgebraicCombinatorics.DominoTilings.rectangle_tileable_iff_even}
\leanhelper
\leanok
If $R_{n,3}$ admits a domino tiling, then $n$ is even (or $n = 0$).
\end{lemma}

\begin{proof}
\leanok
$R_{n,3}$ has $3n$ cells, and each domino covers exactly $2$ cells.
Hence $3n$ must be even. Since $3$ is odd, $n$ must be even.
\end{proof}

\begin{lemma}
\label{lem.details.domino.vertical-trichotomy}
\lean{AlgebraicCombinatorics.DominoTilings.vertical_in_col1_trichotomy}
\leanhelper
\leanok
If a tiling $T$ of $R_{n,3}$ has a vertical domino in column $1$,
then either $T$ has a vertical domino in the \textbf{top} two squares
of column~$1$ (rows $2$--$3$), or $T$ has a vertical domino in the
\textbf{bottom} two squares of column~$1$ (rows $1$--$2$).
\end{lemma}

\begin{proof}
\leanok
A vertical domino in column $1$ covers two adjacent rows.
With only $3$ rows available, the two adjacent rows must be
$\{1,2\}$ or $\{2,3\}$.
\end{proof}

\begin{lemma}
\label{lem.details.domino.base-case-basement}
\lean{AlgebraicCombinatorics.DominoTilings.base_case_basement}
\leanhelper
\leanok
Let $T$ be a tiling of $R_{n,3}$ with $n \ge 2$ that has a vertical domino in the
top two squares of column~$1$. Then $T$ contains the basement domino
$\{(1,1),(2,1)\}$.
\end{lemma}

\begin{proof}
\leanok
The vertical domino covers $(1,2)$ and $(1,3)$. The remaining cell
$(1,1)$ in column~$1$ must be covered by some domino.
It cannot be vertical (the cells above are taken).
It cannot extend left (there is no column~$0$).
So it must be horizontal extending right: $\{(1,1),(2,1)\}$.
\end{proof}

\begin{lemma}[Claim 1]
\label{lem.details.domino.claim1}
\lean{AlgebraicCombinatorics.DominoTilings.claim1_basement_middle_top}
\leanhelper
\leanok
Let $T$ be a faultfree tiling of $R_{n,3}$ that has a vertical domino in the
top two squares of column~$1$.
For each positive integer $i < n/2$, the tiling $T$ contains:
\begin{itemize}
\item the basement domino $\left\{(2i-1,\, 1),\, (2i,\, 1)\right\}$,
\item the middle domino $\left\{(2i,\, 2),\, (2i+1,\, 2)\right\}$,
\item the top domino $\left\{(2i,\, 3),\, (2i+1,\, 3)\right\}$.
\end{itemize}
\end{lemma}

\begin{proof}
By strong induction on $i$.

\emph{Base case ($i = 1$):}
We must prove that $T$
contains the basement domino $\left\{  \left(  1,1\right)  ,\ \left(
2,1\right)  \right\}  $, the middle domino $\left\{  \left(  2,2\right)
,\ \left(  3,2\right)  \right\}  $ and the top domino $\left\{  \left(
2,3\right)  ,\ \left(  3,3\right)  \right\}  $. For the basement domino, we
have already proved this (Lemma~\ref{lem.details.domino.base-case-basement}).
For the middle and the top domino, we argue as
follows: If $T$ had a vertical domino in the $2$-nd column, then this domino
would cover the top two squares of that column (since the bottom square is
already covered by the basement domino $\left\{  \left(  1,1\right)
,\ \left(  2,1\right)  \right\}  $), and thus $T$ would have a fault between
the $2$-nd and $3$-rd columns (since the basement domino $\left\{  \left(
1,1\right)  ,\ \left(  2,1\right)  \right\}  $ also ends at the $2$-nd
column), which would contradict the faultfreeness of $T$. Hence, $T$ has no
vertical domino in the $2$-nd column. Thus, the top two squares of the $2$-nd
column of $T$ must be covered by horizontal dominos. These horizontal dominos
must both protrude into the $3$-rd column (since the corresponding squares in
the $1$-st column are already covered by the left wall), and thus must be the
middle domino $\left\{  \left(  2,2\right)  ,\ \left(  3,2\right)  \right\}  $
and the top domino $\left\{  \left(  2,3\right)  ,\ \left(  3,3\right)
\right\}  $.

\emph{Induction step ($j \to j+1$):}
By the induction hypothesis, $T$ contains the basement domino $\left\{  \left(  2j-1,\ 1\right)
,\ \left(  2j,\ 1\right)  \right\}  $, the middle domino $\left\{  \left(
2j,\ 2\right)  ,\ \left(  2j+1,\ 2\right)  \right\}  $ and the top domino
$\left\{  \left(  2j,\ 3\right)  ,\ \left(  2j+1,\ 3\right)  \right\}  $. We
must now show that $T$ also
contains the basement domino $\left\{  \left(  2j+1,\ 1\right)  ,\ \left(
2j+2,\ 1\right)  \right\}  $, the middle domino $\left\{  \left(
2j+2,\ 2\right)  ,\ \left(  2j+3,\ 2\right)  \right\}  $ and the top domino
$\left\{  \left(  2j+2,\ 3\right)  ,\ \left(  2j+3,\ 3\right)  \right\}  $.

The square $\left(  2j+1,\ 1\right)  $ cannot be covered by a vertical
domino in $T$, since this vertical domino would collide with the middle domino
$\left\{  \left(  2j,\ 2\right)  ,\ \left(  2j+1,\ 2\right)  \right\}  $
(which we already know to belong to $T$). Thus, this square must be covered by
a horizontal domino. This horizontal domino cannot protrude into the $\left(
2j\right)  $-th column (since it would then collide with the basement domino
$\left\{  \left(  2j-1,\ 1\right)  ,\ \left(  2j,\ 1\right)  \right\}  $,
which we already know to belong to $T$), and thus must be the basement domino
$\left\{  \left(  2j+1,\ 1\right)  ,\ \left(  2j+2,\ 1\right)  \right\}  $.

If $T$ had a vertical domino in the $\left(  2j+2\right)  $-nd column, then
this domino would cover the top two squares of that column (since the bottom
square is already covered by the basement domino $\left\{  \left(
2j+1,\ 1\right)  ,\ \left(  2j+2,\ 1\right)  \right\}  $), and thus $T$ would
have a fault between the $\left(  2j+2\right)  $-nd and $\left(  2j+3\right)
$-rd columns (since the basement domino $\left\{  \left(  2j+1,\ 1\right)
,\ \left(  2j+2,\ 1\right)  \right\}  $ also ends at the $\left(  2j+2\right)
$-nd column), which would contradict the faultfreeness of $T$. Hence, $T$ has
no vertical domino in the $\left(  2j+2\right)  $-nd column. Thus, the top two
squares of the $\left(  2j+2\right)  $-nd column of $T$ must be covered by
horizontal dominos. These horizontal dominos must both protrude into the
$\left(  2j+3\right)  $-rd column (since the corresponding squares in the
$\left(  2j+1\right)  $-st column are already covered by the middle domino
$\left\{  \left(  2j,\ 2\right)  ,\ \left(  2j+1,\ 2\right)  \right\}  $ and
the top domino $\left\{  \left(  2j,\ 3\right)  ,\ \left(  2j+1,\ 3\right)
\right\}  $), and thus must be the middle domino $\left\{  \left(
2j+2,\ 2\right)  ,\ \left(  2j+3,\ 2\right)  \right\}  $ and the top domino
$\left\{  \left(  2j+2,\ 3\right)  ,\ \left(  2j+3,\ 3\right)  \right\}  $.
\end{proof}

\begin{lemma}[Claim 2]
\label{lem.details.domino.claim2}
\lean{AlgebraicCombinatorics.DominoTilings.claim2_n_even}
\leanhelper
\leanok
Let $T$ be a faultfree tiling of $R_{n,3}$ that has a vertical domino in the
top two squares of column~$1$. Then $n$ is even.
\end{lemma}

\begin{proof}
\leanok
Since $T$ is a
tiling of $R_{n,3}$ by dominos, the total \# of squares of $R_{n,3}$ must be
even (since each domino covers exactly $2$ squares). But this total \# is
$3n$. Thus, $3n$ must be even, so that $n$ must be even.

(Alternatively, assuming $n$ is odd, one can derive a contradiction
using Claim~1 applied to $i = (n-1)/2$, which forces the cell $(n,1)$
to be covered by a domino that collides with existing ones.)
\end{proof}

\begin{lemma}
\label{lem.details.domino.no-vertical-implies-n-eq-2}
\lean{AlgebraicCombinatorics.DominoTilings.no_vertical_implies_n_eq_2}
\leanhelper
\leanok
If $T$ is a faultfree tiling of $R_{n,3}$ with $n \ge 1$ and no vertical
domino in column~$1$, then $n = 2$.
\end{lemma}

\begin{proof}
\leanok
Since the first column has no vertical domino, each of the three cells
$(1,1)$, $(1,2)$, $(1,3)$ must be covered by a horizontal domino
extending into column~$2$. These three dominos cover all of column~$2$ as
well. If $n > 2$, no domino would span columns $2$ and $3$, creating a
fault at column~$2$, which contradicts faultfreeness.
For $n = 1$, the rectangle $R_{1,3}$ has $3$ cells, which is odd,
so no domino tiling exists.
\end{proof}

\subsection{The classification}

\begin{proposition}
\label{prop.gf.weighted-set.domino.Rn3.ABC}
\lean{AlgebraicCombinatorics.DominoTilings.faultfree_classification}
\leantarget
\leanok
The faultfree domino tilings of
height-$3$ rectangles are precisely the tilings%
\[
A_{2},A_{4},A_{6},A_{8},\ldots,\ \ \ \ \ \ \ \ \ \ B_{2},B_{4},B_{6}%
,B_{8},\ldots,\ \ \ \ \ \ \ \ \ \ C.
\]
More concretely: \medskip

\textbf{(a)} The faultfree domino tilings of a height-$3$ rectangle that
contain a vertical domino in the \textbf{top} two squares of the first column
are $A_{2},A_{4},A_{6},A_{8},\ldots$. \medskip

\textbf{(b)} The faultfree domino tilings of a height-$3$ rectangle that
contain a vertical domino in the \textbf{bottom} two squares of the first
column are $B_{2},B_{4},B_{6},B_{8},\ldots$. \medskip

\textbf{(c)} The only faultfree domino tiling of a height-$3$ rectangle that
contains \textbf{no} vertical domino in the first column is $C$.
\end{proposition}

\begin{proof}
It
clearly suffices to prove parts \textbf{(a)}, \textbf{(b)} and \textbf{(c)}.
We begin with the easiest part, which is \textbf{(c)}: \medskip

\textbf{(c)} Clearly, $C$ is a faultfree domino tiling of $R_{2,3}$ that
contains no vertical domino in the first column. It remains to show that $C$
is the only such tiling.

Indeed, let $T$ be any faultfree domino tiling of $R_{2,3}$ that contains no
vertical domino in the first column. Thus, its first column must be filled
with three horizontal dominos, which all must protrude into the second column
and thus cover that second column as well. If $T$ had any further column, then
$T$ would have a fault between its $2$-nd and its $3$-rd column, which is
impossible for a faultfree tiling. Thus, $T$ must consist only of the three
horizontal dominos already mentioned. In other words, $T=C$. \medskip

\textbf{(a)} It is straightforward to check that the tilings $A_{2}%
,A_{4},A_{6},A_{8},\ldots$ are faultfree (indeed, the basement dominos prevent
faults between the $\left(  2i-1\right)  $-st and $\left(  2i\right)  $-th
columns, whereas the top dominos prevent faults between the $\left(
2i\right)  $-th and $\left(  2i+1\right)  $-st columns). Thus, they are
faultfree domino tilings of a height-$3$ rectangle that contain a vertical
domino in the \textbf{top} two squares of the first column. It remains to show
that they are the only such tilings.

Indeed, let $T$ be any faultfree domino tiling of a height-$3$ rectangle that
contains a vertical domino in the \textbf{top} two squares of the first
column. Let $n$ be the width of this rectangle (so that the rectangle is
$R_{n,3}$). We shall show that $n$ is even, and that $T=A_{n}$.

We know that $T$ contains a vertical domino in the \textbf{top} two squares of
the first column. In other words, $T$ contains the left wall. The remaining square in the first column is the square $\left(
1,1\right)  $, and it must thus be covered by the basement domino $\left\{
\left(  1,1\right)  ,\ \left(  2,1\right)  \right\}  $ (since no other domino
would fit). Hence, $T$ contains the basement domino $\left\{  \left(
1,1\right)  ,\ \left(  2,1\right)  \right\}  $.

Now, Claim 2 (Lemma~\ref{lem.details.domino.claim2}) shows that $n$ is even, so that $n/2$ is a positive integer.
Furthermore, we know that the tiling $T$ contains the left wall, the first
$n-1$ basement dominos
(by Claim 1, Lemma~\ref{lem.details.domino.claim1}), all the middle dominos
and all the top dominos
(also by Claim 1). This leaves only the four squares $\left(  n-1,\ 1\right)
$, $\left(  n,\ 1\right)  $, $\left(  n,\ 2\right)  $ and $\left(
n,\ 3\right)  $ unaccounted for, but there is only one way to tile them:
namely, with the last basement domino $\left\{  \left(  n-1,\ 1\right)
,\ \left(  n,\ 1\right)  \right\}  $ and the right wall $\left\{  \left(
n,2\right)  ,\ \left(  n,3\right)  \right\}  $. Thus, $T$ must contain all the
basement dominos, all the middle dominos, all the top dominos and both walls
(left and right). Since these dominos cover all the squares of $R_{n,3}$, this
entails that $T$ consists precisely of these dominos. In other words,
$T=A_{n}$. \medskip

\textbf{(b)} Proposition \ref{prop.gf.weighted-set.domino.Rn3.ABC}
\textbf{(b)} is just Proposition \ref{prop.gf.weighted-set.domino.Rn3.ABC}
\textbf{(a)} reflected across the horizontal axis of symmetry of $R_{n,3}$.
More precisely, if $T$ is a faultfree tiling with a bottom vertical in
column~$1$, then its reflection $T'$ is faultfree
(Lemma~\ref{lem.details.domino.reflectTiling3-faultfree}) with a top vertical
in column~$1$ (Lemma~\ref{lem.details.domino.reflectTiling3-top-bottom}).
By part~\textbf{(a)}, $T' = A_n$ for some even $n \ge 2$.
Since $B_n$ is the reflection of $A_n$
(Lemma~\ref{lem.details.domino.TilingB-eq-reflect}) and reflection preserves
equality of tilings
(Lemma~\ref{lem.details.domino.reflectTiling3-preserves-equiv}),
reflecting back gives $T = B_n$.

The overall classification follows by
Lemma~\ref{lem.details.domino.vertical-trichotomy}: any faultfree tiling
of $R_{n,3}$ with $n \ge 1$ falls into exactly one of the three cases
(top vertical, bottom vertical, or no vertical in column~$1$).
\end{proof}

\begin{lemma}
\label{lem.details.domino.classification-part-a}
\lean{AlgebraicCombinatorics.DominoTilings.faultfree_top_vertical_classification}
\leanhelper
\leanok
If $T$ is a faultfree tiling of $R_{n,3}$ with a vertical domino in the top
two squares of column~$1$, then $n$ is even and $T = A_n$.
\end{lemma}

\begin{proof}
\leanok
By Claim~2, $n$ is even and $n \ge 2$.
By Claim~1 and the base case lemma, $T$ contains all the dominos of $A_n$.
Since both $T$ and $A_n$ tile the same rectangle, they have the same number
of dominos, and so $T = A_n$.
\end{proof}

\begin{lemma}
\label{lem.details.domino.classification-part-b}
\lean{AlgebraicCombinatorics.DominoTilings.faultfree_bottom_vertical_classification}
\leanhelper
\leanok
If $T$ is a faultfree tiling of $R_{n,3}$ with a vertical domino in the bottom
two squares of column~$1$, then $n$ is even and $T = B_n$.
\end{lemma}

\begin{proof}
\leanok
Reflect $T$ to obtain $T'$, which is faultfree with a top vertical in column~$1$.
By part~\textbf{(a)}, $T' = A_n$ for some even $n \ge 2$.
Since reflection is an involution, $T$ is the reflection of $T'$, and hence
$T$ equals the reflection of $A_n$, which is $B_n$.
\end{proof}

\begin{lemma}
\label{lem.details.domino.classification-part-c}
\lean{AlgebraicCombinatorics.DominoTilings.faultfree_no_vertical_unique}
\leanhelper
\leanok
If $T$ is a faultfree tiling of $R_{2,3}$ with no vertical domino in column~$1$,
then $T = C$.
\end{lemma}

\begin{proof}
\leanok
In a $2 \times 3$ rectangle with no vertical domino in column~$1$,
every domino must be horizontal (since a vertical domino in column~$2$
would force its column-$1$ neighbor to also be vertical in column~$1$, or
would leave a cell that can only be covered by a vertical domino in
column~$1$). Hence the tiling consists of three horizontal dominos
$\{(1,y),(2,y)\}$ for $y = 1, 2, 3$, which is exactly $C$.
\end{proof}

\subsection{Structural helpers for the classification}

\begin{lemma}
\label{lem.details.domino.hasTopVertical-implies-n-ge}
\lean{AlgebraicCombinatorics.DominoTilings.hasTopVerticalInCol_implies_n_ge}
\leanhelper
\leanok
If a tiling $T$ of $R_{n,3}$ has a vertical domino in the top two squares
of some column $c$, then $n \ge c$.
\end{lemma}

\begin{proof}
\leanok
The vertical domino has column coordinate $c$. Since all cells of a
tiling lie in $R_{n,3}$, the column coordinate satisfies $c \le n$.
\end{proof}

\begin{lemma}
\label{lem.details.domino.faultfree-hasTopVertical-implies-n-ge-two}
\lean{AlgebraicCombinatorics.DominoTilings.faultfree_hasTopVerticalInCol_implies_n_ge_two}
\leanhelper
\leanok
If $T$ is a faultfree tiling of $R_{n,3}$ with a vertical domino in the
top two squares of column~$1$, then $n \ge 2$.
\end{lemma}

\begin{proof}
\leanok
By Claim~2 (Lemma~\ref{lem.details.domino.claim2}), $n$ is even.
Since $T$ is nonempty (faultfree implies nonempty), $n \ge 1$,
and being even forces $n \ge 2$.
\end{proof}

\subsection{Height-$2$ Fibonacci recurrence}

These lemmas establish the Fibonacci recurrence $d_{n+2,2} = d_{n,2} + d_{n+1,2}$
using the natural-number-based domino tiling formalization.

\begin{lemma}
\label{lem.details.domino.exists-domino-covering-11}
\lean{AlgebraicCombinatorics.DominoTilings.DominoTiling.exists_domino_covering_11}
\leanhelper
\leanok
For any tiling $T$ of $R_{n+1,2}$ (with $n + 1 \ge 1$), there exists a domino
$d \in T$ that covers the cell $(1,1)$.
\end{lemma}

\begin{proof}
The cell $(1,1) \in R_{n+1,2}$, and by the coverage condition, some
domino in $T$ must contain it.
\end{proof}

\begin{lemma}
\label{lem.details.domino.domino-covering-11-cells}
\lean{AlgebraicCombinatorics.DominoTilings.DominoTiling.domino_covering_11_cells}
\leanhelper
\leanok
If a domino $d$ in a tiling of $R_{n+2,2}$ contains $(1,1)$, then either
$d$ is the vertical domino $\{(1,1),(1,2)\}$ or
$d$ is the horizontal domino $\{(1,1),(2,1)\}$.
\end{lemma}

\begin{proof}
\leanok
The domino is either horizontal or vertical. If horizontal, its second cell
is $(2,1)$ (cannot be $(0,1)$ since $0 < 1$). If vertical, its second cell
is $(1,2)$ (cannot be $(1,0)$ since $0 < 1$).
\end{proof}

\begin{lemma}
\label{lem.details.domino.exists-horizontal-12}
\lean{AlgebraicCombinatorics.DominoTilings.DominoTiling.exists_domino_covering_12_of_horizontal}
\leanhelper
\leanok
If a tiling $T$ of $R_{n+2,2}$ contains a horizontal domino $d$ with
$(1,1) \in d$, then $T$ also contains a domino $d'$ with
$d' = \{(1,2),(2,2)\}$.
\end{lemma}

\begin{proof}
\leanok
The cell $(1,2)$ must be covered by some domino $d'$. This $d'$ cannot
overlap with $d$ (which occupies $(1,1)$ and $(2,1)$). The only
possibility for $d'$ is $\{(1,2),(2,2)\}$.
\end{proof}

\begin{lemma}
\label{lem.details.domino.hasVerticalFirstColumn-or-horizontalPair}
\lean{AlgebraicCombinatorics.DominoTilings.DominoTiling.hasVerticalFirstColumn_or_horizontalPair}
\leanhelper
\leanok
Every tiling $T$ of $R_{n+2,2}$ either has a ``vertical first column''
(a domino with cells $\{(1,1),(1,2)\}$) or has a ``horizontal pair''
(two dominos with cells $\{(1,1),(2,1)\}$ and $\{(1,2),(2,2)\}$).
\end{lemma}

\begin{proof}
\leanok
By Lemma~\ref{lem.details.domino.exists-domino-covering-11}, some domino
covers $(1,1)$. By Lemma~\ref{lem.details.domino.domino-covering-11-cells},
it is either vertical (giving the first case) or horizontal (giving the
second case by Lemma~\ref{lem.details.domino.exists-horizontal-12}).
\end{proof}

\begin{lemma}
\label{lem.details.domino.prependVertical}
\lean{AlgebraicCombinatorics.DominoTilings.DominoTiling.prependVertical}
\leanhelper
\leanok
Given a tiling $T$ of $R_{n,2}$, prepending the vertical domino
$\{(1,1),(1,2)\}$ and shifting all of $T$'s dominos right by $1$ gives a
valid tiling of $R_{n+1,2}$.
\end{lemma}

\begin{proof}
\leanok
The prepended vertical covers column~$1$; the shifted dominos cover
columns $2$ through $n+1$. Pairwise disjointness holds by coordinate
separation.
\end{proof}

\begin{lemma}
\label{lem.details.domino.prependHorizontalPair}
\lean{AlgebraicCombinatorics.DominoTilings.DominoTiling.prependHorizontalPair}
\leanhelper
\leanok
Given a tiling $T$ of $R_{n,2}$, prepending the two horizontal
dominos $\{(1,1),(2,1)\}$ and $\{(1,2),(2,2)\}$ and shifting all of $T$'s
dominos right by $2$ gives a valid tiling of $R_{n+2,2}$.
\end{lemma}

\begin{proof}
\leanok
The two horizontal dominos cover columns~$1$--$2$; the shifted dominos
cover columns $3$ through $n+2$.
\end{proof}

\begin{lemma}
\label{lem.details.domino.restrictAfterVertical}
\lean{AlgebraicCombinatorics.DominoTilings.DominoTiling.restrictAfterVertical}
\leanhelper
\leanok
Given a tiling $T$ of $R_{n+1,2}$ that contains the vertical domino
$\{(1,1),(1,2)\}$, removing it and shifting all other dominos left by $1$
gives a valid tiling of $R_{n,2}$.
\end{lemma}

\begin{proof}
\leanok
After removing the vertical domino, the remaining dominos all have column
coordinates $\ge 2$. Shifting left by $1$ maps them to $R_{n,2}$.
Disjointness and coverage are preserved.
\end{proof}

\begin{lemma}
\label{lem.details.domino.restrictAfterHorizontalPair}
\lean{AlgebraicCombinatorics.DominoTilings.DominoTiling.restrictAfterHorizontalPair}
\leanhelper
\leanok
Given a tiling $T$ of $R_{n+2,2}$ that contains both horizontal
dominos $\{(1,1),(2,1)\}$ and $\{(1,2),(2,2)\}$, removing them and shifting
all other dominos left by $2$ gives a valid tiling of $R_{n,2}$.
\end{lemma}

\begin{proof}
\leanok
After removing the two horizontal dominos, the remaining dominos all have
column coordinates $\ge 3$. Shifting left by $2$ maps them to $R_{n,2}$.
\end{proof}

\begin{lemma}
\label{lem.details.domino.tiling-2-roundtrip}
\lean{AlgebraicCombinatorics.DominoTilings.DominoTiling.tiling_2_roundtrip_TilingEquiv}
\leanhelper
\leanok
For every tiling $T$ of $R_{n+2,2}$, classifying $T$ by
its first-column structure (vertical or horizontal pair), restricting,
and then prepending recovers $T$.
\end{lemma}

\begin{proof}
If $T$ has a vertical first column, the restrict-then-prepend roundtrip
produces a tiling whose dominos are the vertical domino plus all remaining
dominos shifted right then left, recovering the original. The horizontal
pair case is analogous.
\end{proof}

\begin{lemma}
\label{lem.details.domino.tiling-2-from-sum-surjective}
\lean{AlgebraicCombinatorics.DominoTilings.DominoTiling.tiling_2_from_sum_surjective_up_to_TilingEquiv}
\leanhelper
\leanok
The map from $\mathrm{Tiling}(R_{n,2}) \sqcup \mathrm{Tiling}(R_{n+1,2})$
to $\mathrm{Tiling}(R_{n+2,2})$ given by prepending the appropriate dominos
is surjective:
for every tiling $T$ of $R_{n+2,2}$, there exists an element $x$ of
$\mathrm{Tiling}(R_{n,2}) \sqcup \mathrm{Tiling}(R_{n+1,2})$ whose
prepended image equals $T$.
\end{lemma}

\begin{proof}
\leanok
Take $x$ to be the classification of $T$ (by its first-column structure). Then
the prepended image equals $T$ by
Lemma~\ref{lem.details.domino.tiling-2-roundtrip}.
\end{proof}

\begin{lemma}
\label{lem.details.domino.tiling-2-to-sum-from-sum}
\lean{AlgebraicCombinatorics.DominoTilings.DominoTiling.tiling_2_to_sum_from_sum}
\leanhelper
\leanok
For every $x \in \mathrm{Tiling}(R_{n,2}) \sqcup \mathrm{Tiling}(R_{n+1,2})$,
classifying the prepended tiling by its first-column structure and then
restricting recovers $x$.
This shows that the prepending map is injective and
the classification map is surjective.
\end{lemma}

\begin{proof}
\leanok
If $x$ corresponds to a tiling $T$ of $R_{n+1,2}$: prepending a vertical
domino produces a tiling whose first column is vertical,
so the classification recovers $T$ after restriction.
The key is that restricting after prepending a vertical domino recovers $T$.

If $x$ corresponds to a tiling $T$ of $R_{n,2}$: prepending a horizontal pair
produces a tiling whose first column is not vertical, so the classification
recovers $T$ after restriction.
\end{proof}

\subsection{Bridge between $\mathbb{N}$-based and $\mathbb{Z}$-based domino tilings}

\begin{definition}
\label{def.details.domino.bridge-conversions}
\lean{AlgebraicCombinatorics.DominoBridge.toDominoZ, AlgebraicCombinatorics.DominoBridge.toDominoN, AlgebraicCombinatorics.DominoBridge.toTilingZ}
\leanhelper
\leanok
\textbf{(a)} There is a conversion from natural-number-based dominos
(with natural number coordinates) to integer-based dominos (with integer
coordinates) given by casting each coordinate. \medskip

\textbf{(b)} There is a conversion from integer-based dominos
with positive coordinates back to natural-number-based dominos. \medskip

\textbf{(c)} There is a conversion from natural-number-based tilings
of $R_{n,m}$ to integer-based tilings of $R_{n,m}$ obtained by converting
each domino.
\end{definition}

\begin{proof}
\leanok
For the tiling conversion: the converted dominos are pairwise disjoint
(since the coordinate casting preserves disjointness) and their union covers
the integer-based rectangle (by the rectangle correspondence).
\end{proof}

\begin{lemma}
\label{lem.details.domino.rectangle-correspondence}
\lean{AlgebraicCombinatorics.DominoBridge.rectangle_correspondence}
\leanhelper
\leanok
The $\mathbb{N}$-based rectangle $\{(i,j) \mid 1 \le i \le n, 1 \le j \le m\}$ and
the $\mathbb{Z}$-based rectangle $\{(i,j) \in \mathbb{Z}^2 \mid 1 \le i \le n, 1 \le j \le m\}$
are in bijection via the coordinate casting map.
\end{lemma}

\begin{proof}
\leanok
The map is injective (since $\mathbb{N} \hookrightarrow \mathbb{Z}$ is injective)
and surjective onto elements with positive coordinates.
\end{proof}

\begin{lemma}
\label{lem.details.domino.toDominoZ-pairwise-disjoint}
\lean{AlgebraicCombinatorics.DominoBridge.toDominoZ_pairwise_disjoint}
\leanhelper
\leanok
If a set of natural-number-based dominos is pairwise disjoint, then their images under
the coordinate casting are also pairwise disjoint.
\end{lemma}

\begin{proof}
\leanok
Since the coordinate casting preserves the cell set (up to the injective
$\mathbb{N} \hookrightarrow \mathbb{Z}$ embedding), disjointness is preserved.
\end{proof}

\begin{lemma}
\label{lem.details.domino.toDominoZ-cover}
\lean{AlgebraicCombinatorics.DominoBridge.toDominoZ_cover}
\leanhelper
\leanok
If a set of natural-number-based dominos covers the natural-number-based rectangle $R_{n,m}$,
then their images under the coordinate casting cover the $\mathbb{Z}$-based rectangle
$R_{n,m}$.
\end{lemma}

\begin{proof}
\leanok
Every cell $(i,j)$ in the $\mathbb{Z}$-based rectangle corresponds (via
Lemma~\ref{lem.details.domino.rectangle-correspondence}) to a cell in the
natural-number-based rectangle, which is covered by some domino. The image
of that domino under the coordinate casting covers the original cell.
\end{proof}
